
\section{Empirical Validation}
\label{sec:experiments}

The HiPPO dynamics are simple recurrences that can be easily incorporated into various models.
We validate three claims that suggest that when incorporated into a simple RNN, these methods--especially HiPPO-LegS--yield a recurrent architecture with improved memory capability.
In Section~\ref{subsec:membenchmark}, the HiPPO-LegS RNN outperforms
other RNN approaches in benchmark long-term dependency tasks for RNNs.
Section~\ref{subsec:exp-timescale} shows that HiPPO-LegS RNN is much more robust to
timescale shifts compared to other RNN and neural ODE models.
Section~\ref{subsec:exp-scalability} validates the distinct
theoretical advantages of the HiPPO-LegS memory mechanism, allowing fast and
accurate online function reconstruction over millions of time steps.
Experiment details and additional results are described in \cref{sec:experiment-details}.



\paragraph{Model Architectures.}
We first describe briefly how HiPPO memory updates can be incorporated into a
simple neural network architecture, yielding a simple RNN model reminiscent of
the classic LSTM.
Given inputs $x_t$ or features thereof $f_t = u(x_t)$ in any model, the HiPPO
framework can be used to memorize the history of features $f_t$.
Thus, given any RNN update function $h_{t} = \tau(h_{t-1}, x_t)$, we simply replace
$h_{t-1}$ with a projected version of the entire history of $h$, as described in
Figure~\ref{fig:cell}.
The output of each cell is $h_t$, which can be passed through any downstream module (e.g.\ a classification head trained with cross-entropy) to produce predictions.

We map the vector $h_{t-1}$ to 1D with a learned encoding before passing to
$\hippo$ (full architecture in App.~\ref{subsec:model_architectures}).


\subsection{Long-range Memory Benchmark Tasks}
\label{subsec:membenchmark}


\paragraph{Models and Baselines.}
We consider all of the HiPPO methods (\textbf{LegT}, \textbf{LagT}, and \textbf{LegS}).
As we show that many different update dynamics seem to lead to LTI systems that give sensible results (\cref{subsec:high_order_projection}), we additionally consider the \textbf{Rand} baseline that uses random $A$ and $B$ matrices (normalized appropriately) in its updates,
to confirm that the precise derived dynamics are important.
LegT additionally considers an additional hyperparameter $\theta$, which should be set to the timescale of the data if known a priori; to show the effect of the timescale, we set it to the ideal value as well as values that are too large and small.
The \textbf{MGU} is a minimal gated architecture, equivalent to a GRU without the reset gate. The HiPPO architecture we use is simply the MGU with an additional $\hippo$ intermediate layer.

We also compare to several RNN baselines designed for long-term dependencies, including the
\textbf{LSTM}~\cite{lstm}, \textbf{GRU}~\cite{chung2014empirical}, \textbf{expRNN}~\citep{lezcano2019cheap}, and \textbf{LMU}~\citep{voelker2019legendre}.%
\footnote{In our experiments, LMU refers to the architecture in~\citep{voelker2019legendre} while LegT uses the one described in \cref{fig:cell}.}

All methods have the same hidden size in our experiments.
In particular, for simplicity and to reduce hyperparameters, HiPPO variants tie the memory size $N$ to the hidden state dimension $d$, so that all methods and baselines have a comparable number of hidden units and parameters.
A more detailed comparison of model architectures is in \cref{subsec:model_architectures}.



\paragraph{Sequential Image Classification on Permuted MNIST.}
The permuted MNIST (pMNIST) task feeds inputs to a model pixel-by-pixel in the order of a fixed permutation.
The model must process the entire image sequentially -- with non-local structure -- before
outputting a classification label, requiring learning long-term dependencies.

\cref{tab:pmnist} shows the validation accuracy on the pMNIST task for
the instantiations of our framework and baselines.
We highlight that LegS has the best performance of all models.
While LegT is close at the optimal hyperparameter $\theta$, its performance can fall off drastically for a mis-specified window length.
LagT also performs well at its best hyperparameter $\Delta t$.

\cref{tab:pmnist} also compares test accuracy of our methods against reported results from the literature, where the LMU was the state-of-the-art for recurrent models.
In addition to RNN-based baselines, other sequence models have been evaluated on this dataset, despite being against the spirit of the task
because they have global receptive field instead of being strictly sequential.
With a test accuracy of 98.3\%, HiPPO-LegS sets a true state-of-the-art accuracy on the permuted MNIST dataset.



\begin{minipage}{.38\linewidth}%
  \iftoggle{arxiv}{\vspace*{-0.3em}{\vspace*{0.6em}}}
  \centering
  \includegraphics[width=\linewidth]{figures/rnncell2.pdf}
  \captionof{figure}{HiPPO incorporated into a simple RNN model.
  $\hippo$ is the HiPPO memory operator which projects the history of the $f_t$ features depending on the chosen measure.
  }
  \label{fig:cell}
\end{minipage}%
\hfill
\begin{minipage}{.58\linewidth}
    \small
        \centering
        \begin{tabular}{@{}ll@{}}
            \toprule
            Method                     & Val. acc. (\%) \\
            \midrule
            \textbf{-LegS}        & \textbf{98.34}           \\
            -LagT            & 98.15                    \\
            -LegT $\theta = 200$  & 98.0                     \\
            -LegT $\theta = 20$   & 91.75                    \\
            -Rand                 & 69.93 \\
            \midrule
            LMU                        & 97.08                    \\
            ExpRNN                     & 94.67                    \\
            GRU                        & 93.04                    \\
            MGU                        & 89.37                    \\
            RNN                        & 52.98                    \\
            \bottomrule
        \end{tabular}%
        \hfill
        \begin{tabular}{@{}lll@{}}
            \toprule
            Model                                    & Test acc.        \\
            \midrule
            \textbf{HiPPO-LegS}                                & \textbf{98.3} \\
            \midrule
            LSTM~\citep{gu2020improving}              & 95.11         \\
            r-LSTM~\citep{trinh2018learning}         & 95.2          \\
            Dilated RNN~\citep{chang2017dilated}      & 96.1          \\
            IndRNN~\citep{indrnn}                     & 96.0          \\
            URLSTM~\citep{gu2020improving}            & 96.96         \\
            LMU~\citep{voelker2019legendre}           & 97.15         \\
            \midrule
            Transformer~\citep{trinh2018learning}     & 97.9          \\
            TCN~\citep{bai2018empirical} & 97.2          \\
            TrellisNet~\citep{trellisnet}             & 98.13         \\
            \bottomrule
        \end{tabular}
    \captionof{table}{\textbf{(Left)} pMNIST validation, average over 3 seeds. Top: Our methods. Bottom: RNN baselines.
      \textbf{(Right)} Reported test accuracies from previous works. Top: Our
      methods. Middle: Recurrent models. Bottom: Non-recurrent models requiring global receptive field.
    }
    \label{tab:pmnist}
\end{minipage}

\paragraph{Copying task.}
This standard RNN task \cite{arjovsky2016unitary} directly tests memorization, where models must regurgitate a sequence of tokens seen at the beginning of the sequence.
It is well-known that standard models such as LSTMs struggle to solve this task.
\cref{sec:experiment-details}  shows the loss for the Copying task with length $L=200$.
Our proposed update LegS solves the task almost perfectly, while LegT
is very sensitive to the window length hyperparameter.
As expected, most baselines make little progress.








\iftoggle{arxiv}{}{\vspace*{-1em}}

\subsection{Timescale Robustness of HiPPO-LegS}
\label{subsec:exp-timescale}

\paragraph{Timescale priors.}
Sequence models generally benefit from priors on the timescale,
which take the form of additional hyperparameters in standard models.
Examples include the ``forget bias'' of LSTMs which needs to be modified to address long-term dependencies~\cite{jozefowicz2015empirical,tallec2018can},
or the discretization step size $\Delta t$ of HiPPO-Lag and HiPPO-LegT (\cref{sec:discretization}).
The experiments in \cref{subsec:membenchmark} confirm their importance.
\cref{fig:copy200} (Appendix) and \cref{tab:pmnist} ablate these hyperparameters,
showing that for example the sliding window length $\theta$ must be set correctly for LegT.
Additional ablations for other hyperparameters are in \cref{sec:experiment-details}.



\paragraph{Distribution shift in trajectory classification.}
Recent trends in ML have stressed the importance of understanding robustness under distribution shift, when training and testing distributions are not i.i.d.
For time series data, for example, models may be trained on EEG data from one hospital, but deployed at another using instruments with different sampling rates~\citep{shah2018temple, saab2020weak}; or a time series may involve the same trajectory evolving at different speeds.
Following~\citet{kidger2020neural}, we consider the Character Trajectories dataset~\citep{bagnall2018uea},
where the goal is to classify a character from a sequence of pen stroke measurements, collected from one user at a fixed sampling rate.
To emulate timescale shift (e.g.\ testing on another user with slower handwriting),
we consider two standard time series generation processes:
(1) In the setting of sampling an underlying sequence at a fixed rate, we change the test sampling rate; crucially, the sequences are variable length so the models are unable to detect the sampling rate of the data.
(2) In the setting of irregular-sampled (or missing) data with timestamps, we scale the test timestamps.





Recall that the HiPPO framework models the underlying data as a continuous function and interacts with discrete input only through the discretization.
Thus, it seamlessly handles missing or irregularly-sampled data by simply evolving according to the given discretization step sizes (details in \cref{sec:discretization-full}).
Combined with LegS timescale invariance (Prop.~\ref{prop:timescale}), we expect HiPPO-LegS to work automatically in all these settings.
We note that the setting of missing data is a topic of independent interest and we compare against SOTA methods, including the GRU-D~\cite{che2018recurrent} which learns a decay between observations,
and neural ODE methods which models segments between observations with an ODE.

\cref{tab:charactertrajectories} validates that standard models can go catastrophically wrong when tested on sequences at different timescales than expected.
Though all methods achieve near-perfect accuracy ($\geq$ 95\%) without distribution shift,
aside from HiPPO-LegS, no method is able to generalize to unseen timescales.




\begin{table}[ht]
    \small
    \centering
    \caption{Test set accuracy on Character Trajectory classification on out-of-distribution timescales.   }
    \begin{tabular}{lccccccc}
        \toprule
        Model                     & \textbf{LSTM} & \textbf{GRU} & \textbf{GRU-D} & \textbf{ODE-RNN} & \textbf{NCDE} & \textbf{LMU} & \textbf{HiPPO-LegS} \\
        \midrule
        100Hz $\to$ 200Hz         & 31.9          & 25.4         & 23.1           & 41.8             & 44.7          & 6.0          & \textbf{88.8}       \\
        200Hz $\to$ 100Hz         & 28.2          & 64.6         & 25.5           & 31.5             & 11.3          & 13.1         & \textbf{90.1}       \\
        \midrule
        Missing values upsample   & 24.4          & 28.2         & 5.5            & 4.3              & 63.9          & 39.3         & \textbf{94.5}       \\
        Missing values downsample & 34.9          & 27.3         & 7.7            & 7.7              & 69.7          & 67.8         & \textbf{94.9}       \\
        \bottomrule
    \end{tabular}
    \label{tab:charactertrajectories}
\end{table}




\subsection{Theoretical Validation and Scalability}
\label{subsec:exp-scalability}

We empirically show that HiPPO-LegS can scale to capture dependencies
across millions of time steps, and its memory updates are computationally
efficient (processing up to 470,000 time steps/s).

\paragraph{Long-range function approximation.}
We test the ability of
different memory mechanisms in approximating an input function, as described in
the problem setup in Section~\ref{subsec:hippo-setup}.
The model only consists of the memory update (Section~\ref{sec:theory_legs}) and
not the additional RNN architecture.
We choose random samples from a continuous-time band-limited white noise
process, with length $10^6$.
The model is to traverse the input sequence, and then asked to reconstruct the
input, while maintaining no more than 256 units in memory (\cref{fig:function_approx}).
This is a difficult task; the LSTM fails with even sequences of length 1000 (MSE
$\approx$ 0.25).
As shown in Table~\ref{tab:mse_speed}, both the LMU and HiPPO-LegS are able to
accurately reconstruct the input function, validating that HiPPO can solve the
function approximation problem even for very long sequences.
\cref{fig:function_approx} illustrates the function and its approximations, with
HiPPO-LegS almost matching the input function while LSTM unable to do so.

\paragraph{Speed.}
HiPPO-LegS operator is computationally efficient both in theory
(Section~\ref{sec:theory_legs}) and in practice.
We implement the fast update in C++ with Pytorch binding and show in
Table~\ref{tab:mse_speed} that it can perform 470,000 time step updates per second on a single CPU
core, 10x faster than the LSTM and LMU.\footnote{The LMU is only known to be fast with the simple forward Euler
discretization~\citep{voelker2019legendre}, but not with more sophisticated methods such as bilinear and
ZOH that are required to reduce numerical errors for this task.}

\begin{minipage}{.4\linewidth}
  \small
  \centering
  \begin{tabular}{lll}
      \toprule
      Method     & Error  & Speed  \\
      & (MSE) & (elements / sec) \\
      \midrule
      LSTM       & 0.25 & 35,000             \\
      LMU        & 0.05 & 41,000             \\
      HiPPO-LegS & 0.02 & 470,000            \\
      \bottomrule
  \end{tabular}
  \captionof{table}{Function approximation error after 1 million time steps, with 256 hidden units.}
  \label{tab:mse_speed}
\end{minipage}%
\hfill
\begin{minipage}{.6\linewidth}%
  \centering
  \includegraphics[width=\linewidth]{figures/function_approx.pdf}
  \captionof{figure}{Input function and its reconstructions.}
  \label{fig:function_approx}
\end{minipage}%

\subsection{Additional Experiments}
\label{subsec:sequencetasks}

We validate that the HiPPO memory updates also perform well on more generic
sequence prediction tasks not exclusively focused on memory.
Full results and details for these tasks are in \cref{sec:experiment-details}.

\paragraph{Sentiment classification task on the IMDB movie review dataset.}
Our RNNs with HiPPO memory updates
perform on par with the LSTM, while other 
long-range memory approaches such as expRNN perform poorly on this more generic task
(\cref{sec:imdb}).


\paragraph{Mackey spin glass prediction.} This physical simulation task tests the
ability to model chaotic dynamical systems. HiPPO-LegS outperforms the LSTM,
LMU, and the best hybrid LSTM+LMU model from~\cite{voelker2019legendre},
reducing normalized MSE by $30\%$ (\cref{sec:mackey}).

